% Written by projects/explain-top-move/paper/generate.mjs; do not edit by hand.
\begin{minipage}[t]{0.4\textwidth}\vspace{0pt}
\chessboard[setfen={r1b2rk1/p1p1qpbp/1pn3p1/3p4/N2P4/1Q2PPN1/PP4PP/R3KB1R b KQ - 1 14}, boardfontsize=13pt, pgfstyle=straightmove, arrow=to, linewidth=0.07em, color=red!75!black, markmoves={g7-d4}, pgfstyle=straightmove, arrow=to, linewidth=0.07em, color=green!50!black, markmoves={c6-d4}]
\end{minipage}\hfill
\begin{minipage}[t]{0.56\textwidth}\vspace{0pt}\small
\textbf{Played:} 14\ldots{}\symbishop{}xd4, Black to move; the app's verdict: Inaccuracy.\par
\textbf{Engine:} +5.0 with the best move, +3.5 after yours.\par
\textbf{Best line:} 14\ldots{}\symknight{}xd4 15.~\symqueen{}d3 \symbishop{}f5 16.~\symqueen{}d2 \symknight{}c2+ 17.~\symking{}f2\par
\textbf{UCI:} \uci{g7d4}; best line \uci{c6d4 b3d3 c8f5 d3d2 d4c2 e1f2}\par
\textbf{FEN:} \fen{r1b2rk1/p1p1qpbp/1pn3p1/3p4/N2P4/1Q2PPN1/PP4PP/R3KB1R b KQ - 1 14}
\end{minipage}

\smallskip\small
\begin{tabularx}{\textwidth}{>{\raggedright\arraybackslash}p{0.17\textwidth} L L}
\toprule
section & before & after \\
\midrule
What Nxd4 does instead & it attacks the queen on b3 and the pawn on f3; it wins the pawn outright; if they move the queen (Qxd5), Qxe3+ Be2 Nc2+ Kf1 Bd4 Ne4 Bf5 Nac3 Nxa1, and you have won about 6 pawns of material \textperiodcentered{} 14...Nxd4 15.Qd3 Bf5 16.Qd2 Nc2+ 17.Kf2 Nxa1 18.Nxf5 & it attacks the queen on b3 and the pawn on f3; it wins the pawn outright; if they move the queen (Qxd5), you have won about 6 pawns of material by move 19 \textperiodcentered{} 14...Nxd4 15.Qd3 Bf5 16.Qd2 Nc2+ 17.Kf2 Nxa1 18.Nxf5 \\
\bottomrule
\end{tabularx}

\textbf{The judge (run 1, where the new text was X):} ``Y shows the forcing sequence after Qxd5 (Qxe3+, Nc2+ and finally Nxa1) that wins the rook, while X only says it happens 'by move 19'.''
